%% toc: P2 15 15.5,15.6,15.7
\section{Horn closure as a least fixed point, and the proof route for Cook--Levin}
\label{sec:p02-w02b-p02-015-links}

This section makes two connections that the satisfiability procedures of
\cref{sec:p02-w02-p02-019-sat-procedures} need beyond Part~2.  The first
identifies the forward closure of a Horn list
(Definition~\ref{def:p02-w02-p02-019-forward-closure}) with the least fixed
point \(M_P\) of the immediate-consequence operator \(T_P\) of Part~9
(Definition~\ref{def:p09-definite-program}), for the propositional definite
program read off the list.  Its inputs are a Horn list and its definite
clauses; its conclusions are the equality of the two least objects, the least
Herbrand model property for the variable-free case, and the Horn-SAT test
restated as a goal test against \(M_P\).  The second supplies the written
proof of the Cook--Levin Theorem~\ref{thm:p04-cook-levin} under the machine,
reduction and class conventions of Part~4
(Convention~\ref{def:p04-complexity-convention},
Definition~\ref{def:p04-resource-classes}).  Its inputs are those conventions
and the Part~2 results on certificates and tableau formulas; its conclusion is
that \(\mathsf{SAT}\) is \(\mathsf{NP}\)-complete in the sense fixed in
Part~4.

\subsection{Horn lists as propositional definite programs}
\label{sec:p02-w02b-p02-015-horn-programs}

\begin{hbdefinition}{def:p02-w02b-p02-015-horn-program}{The definite program of a Horn list}
Let \(S\) be a Horn list (Definition~\ref{def:p02-tractable-classes}) with
variable set \(A_S\), and use the body \(B(C)\), the head
\(\operatorname{hd}(C)\) and the split into definite and negative clauses of
Definition~\ref{def:p02-w02-p02-019-horn-parts}.  Read every variable of
\(A_S\) as a nullary relation symbol of a first-order signature whose relation
symbols are exactly the members of \(A_S\).  By
Definition~\ref{def:p02-herbrand-universe-base} a nullary relation symbol
contributes its own atom and no other, so the Herbrand base over this
signature is \(A_S\).  The \emph{definite program of \(S\)}, written
\(P_S\), contains, for each definite clause \(C\) of \(S\) with
\(B(C)=\{q_1,\ldots,q_k\}\), listed in increasing order of index, and
\(\operatorname{hd}(C)=r\), the definite clause
\[
 r\leftarrow q_1,\ldots,q_k
\]
of Definition~\ref{def:p09-definite-program}; a positive unit clause
\(\{r\}\) gives the clause \(r\leftarrow{}\) with empty body.  A negative
clause \(C\) of \(S\) with \(B(C)=\{q_1,\ldots,q_k\}\) gives no program
clause.  It is read as the \emph{goal} \(\leftarrow q_1,\ldots,q_k\), which is
violated exactly when all of \(q_1,\ldots,q_k\) are true; the empty clause is
the goal with empty body, violated by every valuation.

No clause of \(P_S\) contains a first-order variable, so the empty
substitution is the only ground substitution that matters, and adding a fresh
constant under the convention of Definition~\ref{def:p09-definite-program}
adds no atom to the base.  Accordingly \(\mathit{HB}_{P_S}=A_S\), and
\(T_{P_S}\) and \(M_{P_S}\) are the operator and least fixed point of that
definition for the program \(P_S\).  For a set \(I\subseteq A_S\), \(v_I\) is
the valuation of Definition~\ref{def:p02-w02-p02-019-horn-parts}, which makes
exactly the members of \(I\) true; it is the propositional form of the
Herbrand interpretation \(I\).
\index{definite program!of a Horn list}\index{Horn clause!as program clause}
\end{hbdefinition}

\begin{hbtheorem}{thm:p02-w02b-p02-015-horn-lfp}{proved=annex:p02-w02b-p02-015-horn-lfp}{Forward closure is the least model of the definite program}
Let \(S\) be a Horn list, \(T=T_{P_S}\) on subsets of \(A_S\), and
\(F=\operatorname{Fwd}(S)\).
\begin{enumerate}
  \item For every \(I\subseteq A_S\),
        \(T(I)=\{\operatorname{hd}(C): C\text{ a definite clause of }S,\
        B(C)\subseteq I\}\), and \(T\) is monotone.
  \item A set \(X\subseteq A_S\) is closed under \(S\) iff
        \(T(X)\subseteq X\).
  \item \(T^k(\varnothing)\subseteq T^{k+1}(\varnothing)\) for every
        \(k\in\N\); \(T(F)=F\); \(F\subseteq X\) for every \(X\subseteq A_S\)
        with \(T(X)\subseteq X\); and
        \[
         F=\bigcup_{k<\omega}T^k(\varnothing)=M_{P_S}=T^k(\varnothing)
         \qquad\text{for every }k\geq|A_S|.
        \]
  \item For every \(I\subseteq A_S\), \(v_I\) satisfies every definite clause
        of \(S\) iff \(T(I)\subseteq I\).  Hence \(M_{P_S}\) is the least
        Herbrand model of \(P_S\), and an atom \(a\in A_S\) is true under every
        valuation satisfying all definite clauses of \(S\) iff
        \(a\in M_{P_S}\).
  \item \(S\) is unsatisfiable iff some negative clause \(C\) of \(S\) has
        \(B(C)\subseteq M_{P_S}\), that is, iff the least model of \(P_S\)
        violates one of the goals of \(S\).  If \(S\) is satisfiable,
        \(v_{M_{P_S}}\) satisfies \(S\) and is its least model.
\end{enumerate}
\end{hbtheorem}

Clause~4 is the variable-free case of the least-model
Theorem~\ref{thm:p09-definite-fixpoint}, which Part~9 states with a citation
to \textcite{VanEmdenKowalski1976}; the annex proof above is complete for
programs all of whose relation symbols are nullary.  Clause~3 says that
\(\operatorname{Fwd}(S)\), defined in Part~2 as the least set closed under
the definite clauses, is the same set as the iterate
\(\bigcup_kT_P^k(\varnothing)\) that Part~9 uses to connect model semantics
with Datalog evaluation and SLD resolution
(\cref{sec:p09-logic-programming}).  The two characterizations of the least
model in clause~4 of Theorem~\ref{thm:p02-w02-p02-019-horn-closure} and
clause~5 above are therefore one statement.  By
Theorem~\ref{thm:p02-w02-p02-019-horn-procedure}, the marked set of
Method~\ref{def:p02-w02-p02-019-horn-method} after its propagation step is
\(\operatorname{Fwd}(S)\), so that method computes \(M_{P_S}\) within
\(O(N)\) steps: each atom leaves the queue once, and each body occurrence
lowers one counter once, whereas the iterates \(T^k(\varnothing)\) are each
defined from the whole previous stage.  The order in which atoms leave the
queue does not affect the result, since the result is the least fixed point.

\begin{hbexample}{def:p02-w02b-p02-015-horn-program-example}{Iterating \(T_P\) on the lists of the forward-chaining example}
Take the list \(S=(C_1,C_2,C_3,C_4)\) of
Example~\ref{def:p02-w02-p02-019-horn-example}, with
\(C_1=\{p_1\}\), \(C_2=\{\neg p_1,p_2\}\),
\(C_3=\{\neg p_2,\neg p_4,p_3\}\) and \(C_4=\{\neg p_1,\neg p_3\}\).  Then
\(A_S=\{p_1,p_2,p_3,p_4\}\), the program \(P_S\) is
\[
 p_1\leftarrow{},\qquad p_2\leftarrow p_1,\qquad p_3\leftarrow p_2,p_4,
\]
and \(C_4\) is the goal \(\leftarrow p_1,p_3\).  The iterates are
\(T(\varnothing)=\{p_1\}\), \(T^2(\varnothing)=\{p_1,p_2\}\) and
\(T^3(\varnothing)=\{p_1,p_2\}\): the third clause needs \(p_4\), which is the
head of no clause.  So \(M_{P_S}=\{p_1,p_2\}\), the same set that forward
chaining marked in the example.  The goal body \(\{p_1,p_3\}\) is not
contained in it, so by clause~5 \(S\) is satisfiable, and
\(v_{\{p_1,p_2\}}\) is its least model.

Appending \(C_5=\{p_4\}\) adds the program clause \(p_4\leftarrow{}\).  Now
\(T(\varnothing)=\{p_1,p_4\}\), \(T^2(\varnothing)=\{p_1,p_2,p_4\}\),
\(T^3(\varnothing)=\{p_1,p_2,p_3,p_4\}\), and \(T^4(\varnothing)\) is the same
set, reached within \(|A_S|=4\) rounds as clause~3 guarantees.  The goal body
\(\{p_1,p_3\}\) now lies in the least model, and the extended list is
unsatisfiable, as step~(4) of the method found.
\end{hbexample}

\begin{hbwarning}{def:p02-w02b-p02-015-two-sat-programs}{2-CNF clauses that are not program clauses}
The program reading needs every clause to have at most one positive literal.
The 2-CNF clause \(\{p_1,p_2\}\) has two; it is neither a definite clause nor
a goal, \(T_P\) of Definition~\ref{def:p09-definite-program} is not defined
for it, and the list consisting of it alone has no least model
(Warning~\ref{def:p02-w02-p02-019-horn-limit}).  Satisfiability of such lists
is decided by the implication graph instead
(Theorem~\ref{thm:p02-w02-p02-019-two-sat}).  The valuation \(v_S\) that
Method~\ref{def:p02-w02-p02-019-two-sat-method} extracts is chosen by the
numbering of strong components, not as a least fixed point.  When a 2-CNF list
\(S\) is also Horn and satisfiable, \(v_S\) satisfies \(S\), so by clause~2 of
Theorem~\ref{thm:p02-w02-p02-019-horn-closure} it makes every atom of
\(M_{P_S}\) true; the two procedures then agree on satisfiability, and the
least model is the one that forward chaining returns.
\end{hbwarning}

\subsection{The Part~4 conventions and the proof of the Cook--Levin theorem}
\label{sec:p02-w02b-p02-015-cook-levin}

Theorem~\ref{thm:p04-cook-levin} is stated under the conventions of Part~4:
\(\mathsf{NP}\) is \(\bigcup_k\mathsf{NTIME}(n^k)\) for nondeterministic
machines with a read-only input tape
(Definition~\ref{def:p04-resource-classes}), and reductions are computed by
deterministic machines of that kind
(Convention~\ref{def:p04-complexity-convention}).  The Part~2 results that
carry the combinatorial work use certificates and machines that may write on
their input tape: Convention~\ref{def:p02-complexity-standing},
Convention~\ref{def:p02-w02-p02-018-verifier},
Theorem~\ref{thm:p02-w02-p02-018-membership} and the tableau construction of
Theorem~\ref{thm:p02-w02-p02-018-cook-levin}.  The convention and lemma below
connect the two settings, and the annex proof of
Theorem~\ref{thm:p04-cook-levin} combines them.

\begin{hbconvention}{def:p02-w02b-p02-015-machine-reading}{Reading the Part~4 machines}
A machine of Convention~\ref{def:p04-complexity-convention} is read in the
format fixed at the start of the proof in
Annex~\ref{annex:p02-w02-p02-018-cook-levin}: tapes \(1,\ldots,\kappa\) with
cells \(0,1,2,\ldots\); the input in cells \(0,\ldots,n-1\) of tape~\(1\)
and blanks elsewhere; every head on cell \(0\) at the start; a left move at
cell \(0\) leaving the head there; and halting states \(q_{\mathrm a}\) and
\(q_{\mathrm r}\).  Tape~\(1\) is read-only, that is, the symbol written there
is always the one scanned.  A \emph{nondeterministic} machine replaces
\(\delta\) by a finite relation that gives, for each nonhalting state \(s\)
and tuple \(\vec a\) of scanned symbols, a finite and possibly empty set
\(\Delta(s,\vec a)\) of triples \((s',\vec b,\vec d)\).  A \emph{branch} on
input \(x\) is a finite or infinite sequence of configurations that starts in
the initial configuration, passes from each configuration to one obtained by
a triple of \(\Delta\), and is finite exactly when it ends in a configuration
with a halting state or with an empty \(\Delta(s,\vec a)\); the second kind
of end counts as rejection.  A machine computes a map \(f\) when it halts on
each \(x\) with \(f(x)\) in cells \(0,\ldots,\operatorname{len}(f(x))-1\) of a
designated output tape other than tape~\(1\), followed by blanks.  The same
output format is used for the polynomial-time maps of
Convention~\ref{def:p02-complexity-standing}.  This reading is an assumption
of Lemma~\ref{thm:p02-w02b-p02-015-bridge}; Part~4 does not fix these
details.
\index{nondeterministic machine!branch}
\end{hbconvention}

\begin{hblemma}{thm:p02-w02b-p02-015-bridge}{proved=annex:p02-w02b-p02-015-bridge}{Part~2 and Part~4 conventions for maps, reductions and \(\mathsf{NP}\)}
Read machines as in Convention~\ref{def:p02-w02b-p02-015-machine-reading}.
\begin{enumerate}
  \item A map \(f:\{0,1\}^*\to\{0,1\}^*\) is computable within a polynomial
        number of steps by a deterministic multitape machine whose input tape
        may be written iff it is so computable by one whose input tape is
        read-only.  The same holds for deciding a language.  Hence the
        polynomial-time many-one reductions of
        Convention~\ref{def:p02-complexity-standing} and the relation
        \(\leq_m^p\) of Convention~\ref{def:p04-complexity-convention} are the
        same relation, and the two definitions of \(\mathsf P\) give the same
        class.
  \item If \(L\in\mathsf{NP}\) in the sense of
        Definition~\ref{def:p04-resource-classes}, then \(L\in\mathsf{NP}\)
        in the sense of Convention~\ref{def:p02-complexity-standing}, with
        the predicate read as in
        Convention~\ref{def:p02-w02-p02-018-verifier}.
  \item If a polynomial-time predicate \(R\), in the sense of
        Convention~\ref{def:p02-w02-p02-018-verifier}, satisfies, for every
        bit string \(x\) of length \(n\),
        \[
         x\in L\quad\text{iff}\quad\text{some }y\text{ with }
         \operatorname{len}(y)\leq n\text{ satisfies }R(x,y),
        \]
        then \(L\in\mathsf{NP}\) in the sense of
        Definition~\ref{def:p04-resource-classes}.
\end{enumerate}
\end{hblemma}

Clause~3 is the case of certificate bound \(q(n)=n\), which is the bound of
Theorem~\ref{thm:p02-w02-p02-018-membership}; it is all that the proof of
Theorem~\ref{thm:p04-cook-levin} needs.  The equality of the two classes for
an arbitrary polynomial bound, which also requires the machine to lay out
\(q(n)\) cells before guessing, is \cite[thm.~2.6]{AroraBarak2009}.

\begin{hbexample}{def:p02-w02b-p02-015-choice-code}{Choices coded as a certificate}
Suppose that in some configuration a machine has three applicable triples,
listed in a fixed order as \(O_0,O_1,O_2\), and that no set
\(\Delta(s,\vec a)\) has more than three members.  The proof of clause~2
codes each choice by a block of \(\beta=2\) bits.  The blocks \(00\), \(01\)
and \(10\) select \(O_0\), \(O_1\) and \(O_2\).  The block \(11\) has value
\(3\), and \(3\bmod 3=0\) selects \(O_0\) again.  So every bit string
determines a branch, and every branch of length \(T'\) is determined by a
string of length \(2T'\).  This is the certificate that clause~2 hands to the
tableau construction of Theorem~\ref{thm:p02-w02-p02-018-cook-levin}.
\end{hbexample}

With Lemma~\ref{thm:p02-w02b-p02-015-bridge}, membership of \(\mathsf{SAT}\)
in \(\mathsf{NP}\) under Definition~\ref{def:p04-resource-classes} follows
from the certificate check of
Theorem~\ref{thm:p02-w02-p02-018-membership} by clause~3.  Hardness follows
from clause~2, the tableau formula of
Theorem~\ref{thm:p02-w02-p02-018-cook-levin}, and clause~1.  The language is
the \(\mathsf{SAT}\) of Definition~\ref{def:p02-sat-taut}, which Part~4 uses
(\cref{sec:p04-complexity-complete-problems}); its code table is the one
fixed in Convention~\ref{def:p02-w02-p02-018-input}.  The full argument is
Annex~\ref{annex:p02-w02b-p02-015-cook-levin}.  Part~9 presents the same
tableau method in the form of Cook's paper and records how Cook's
query-machine reducibility differs from the many-one reduction proved here
(\cref{sec:p09-presentation-n-c-026}).

\begin{hbwarning}{def:p02-w02b-p02-015-cook-levin-scope}{What this proof route does and does not establish}
The proof establishes \(\mathsf{NP}\)-completeness of the language of all
satisfiable encoded formulas under polynomial-time many-one reductions.  It
does not prove completeness of \(3\mathsf{SAT}\), which
Table~\ref{tab:p04-completeness-catalog} cites, and it does not prove the
logspace form of the reduction.  It also does not establish Cook's own
statement: \textcite{Cook1971} uses polynomial-time query-machine
reducibility to the tautologies in disjunctive normal form, a different
reduction notion (\cref{sec:p09-presentation-n-c-026}).  The equality of the
two \(\mathsf{NP}\) definitions is used only in the directions of clauses~2
and~3 of Lemma~\ref{thm:p02-w02b-p02-015-bridge}.  The lemma also depends on
the machine format of
Convention~\ref{def:p02-w02b-p02-015-machine-reading}; a different format for
Part~4 machines would need its own simulation argument.
\end{hbwarning}
